% !TEX root = ../PhD Thesis.tex

\chapter*{Preface and Acknowledgements}
\addcontentsline{toc}{chapter}{Acknowledgements}

“Conformity is the death of progress.”
\\
\\
It was this motto that I wielded without ever yielding. I never truly felt like a scientist. I am disorganised, and I am especially sensitive to failure, all in a field where failure is not incidental but endemic. I have always resisted the canonical academic path, in which a scientist devotes an entire career to a single gene, protein, pathway, or at least a narrow thematic corridor. Instead, I always preferred ever-shifting terrain: exploratory projects, changing questions, intellectual migration.
\\
I dislike repetition. I dislike the sensation of stagnant progress. The feeling that one sacrifices mental health for a marginal increment of data, like a World War I general sacrificing men for a few inches of land. The analogy is cold, disproportionate, and I recognise that fully. Yet the very impulse to make it is revealing.
\\
Although I was a strong student overall, my highest marks were always in the humanities: languages and history. For many years, I wanted to be a creative writer, to craft novels with morally ambiguous characters, where the hero confronts a great external conflict only to discover that the true struggle is internal. How much of himself is he willing to sacrifice for his goal? In that regard, life imitated art. My world, too, contains grey characters: stress granules as both tumour suppressors and tumour promoters. And though I am no hero, I have repeatedly faced a similar question. How much of myself should be sacrificed for scientific progress?
\\
The project itself was not without turbulence. At its inception, it was challenged on conceptual grounds: the heterogeneity of stressors would inevitably produce heterogeneous stress granules. If the upstream perturbations vary widely, how could the downstream assemblies be consistent? My preliminary data, precisely the differential centrifugation studies, suggested otherwise. The granules appeared remarkably conserved.
\\
Of course, there is a marked irony in this. The very data that secured the project’s legitimacy are now the data I regard as methodologically inadequate. That does not imply that stress granules induced by distinct stressors share nothing in common. It suggests only that the question is more complex than I once believed. I remain convinced that all knowledge, including “negative” knowledge, as was so called in regards to my project, as invaluable to scientific progress.
\\
As a neurodivergent individual, I often felt out of place, though never entirely without one. I am obsessive. I cannot rest when a question lacks an answer. Many nights were spent carousing with insomnia. I am perhaps more imaginative than most, capable of generating hypotheses that bridge conceptual gaps others might not attempt to cross. Yet imagination is double-edged: when one sees connections everywhere, even the impossible can appear plausible. A hypothesis that felt visionary to me may have seemed quixotic to others.
\\
I have always been competitive, but situated at a crossroads. I know less mathematics than engineers, less molecular biology than career biologists. Mediocre in many domains, but great at none. Watching peers excel in clearly defined strengths, I often felt diffuse by comparison. And so I leaned on other faculties: stubbornness, diligence, an unhealthy disregard for personal cost.
\\
Within that darkness, I often took proverbial shots in the dark. I pursued techniques and ideas others would not entertain: simulating RNA sequencing for marginal informational gain; resurrecting equations older than myself to calculate how a single RNA transcript might sediment under specific centrifugation regimes, assuming, improbably, that it behaved as a sphere influenced only by mass; estimating how two molecular additions, individually negligible yet cumulatively significant, would alter that sedimentation. I explored how specific RNA-binding proteins regulate not only themselves but one another. Much of it bordered on the unreasonable. Most of it I did not particularly enjoy.
\\
My preference has always been different: to use the keys I already possess to open new doors. Yet I was seldom afforded that luxury. Relentlessness became my default mode, less a virtue than a coping mechanism. Failure felt personal, with success merely feeling like a failure that was yet to come. Whenever a question was answered and closed, it simply brought relief, but never celebration.
\\
Perhaps, when I read this again in the future, I will be more committed to the pursuit of doors rather than the compulsive forging of new keys. Or at least, I will believe that the keys I forged have contributed, however modestly, to the "bettering of humankind".
\\
And now, an abrupt change, from the bleak paragraph (albeit less so than what is my personal signature), to the concomitant section of acknowledgments, which I chose to maintain herein. 


DONT FORGET TO COMPLETE THE ACKNOWLEDGMENTS YOU DIMWIT